
\documentclass{beamer}
\usetheme{Madrid}

\title{Effects of Automation and Artificial Intelligence on Employment}
\subtitle{Contemporary Economic Issues: ECON 204}
\author{Muhebullah Karimzada}
\date{Winter 2025}
\begin{document}

\frame{\titlepage}

\begin{frame}{Labor Market Polarization}
Automation and AI tend to disproportionately affect different types of jobs, leading to labor market polarization.
\vspace{0.5cm}\\
\textbf{High-Skilled Jobs}:
\begin{itemize}
    \item  Automation complements tasks that require advanced cognitive, analytical, and creative skills (e.g., software development, data analysis). 
    \item Workers in these areas often see increased demand and higher wages due to AI-enabled productivity boosts.
\end{itemize}
\end{frame}

\begin{frame}{Labor Market Polarization}
Automation and AI tend to disproportionately affect different types of jobs, leading to labor market polarization.
\vspace{0.5cm}\\
 \textbf{Low-Skilled Jobs}: 
 \begin{itemize}
    \item Certain low-skilled, non-automatable jobs (e.g., caregiving, hospitality) may also remain in demand, particularly where human interaction is essential.

\end{itemize}

   \textbf{Middle-Skilled Jobs}: 
    \begin{itemize}
  \item Routine, middle-skill tasks (e.g., clerical work, assembly line operations) are most vulnerable to displacement by AI and robotics, leading to a decline in middle-skill job opportunities.
\end{itemize}

\end{frame}






\begin{frame}{Job Displacement and Job Creation}
\textbf{Job Destruction:}
\begin{itemize}
    \item Automation replaces tasks that are repetitive, routine, and predictable (e.g., manufacturing, bookkeeping, transportation).
    \item A 2017 study by Frey and Osborne estimated that about 47\% of U.S. jobs are at high risk of automation.
\end{itemize}
\vspace{0.5cm}
\textbf{Job Creation:}
\begin{itemize}
    \item AI and automation also give rise to new industries, occupations, and tasks (e.g., data scientists, AI specialists, robotics engineers).
    \item Expansion in sectors like renewable energy, healthcare, and AI-driven tech startups.
\end{itemize}
\end{frame}




\begin{frame}{Job Displacement and Job Creation}

\textbf{Net Effect:}
\begin{itemize}
    \item In the short term, job destruction often outpaces job creation, causing unemployment and worker displacement.
    \item In the long term, complementary industries and tasks may emerge, but retraining and education are critical to ensuring workers can transition to these roles.
\end{itemize}
\end{frame}





\begin{frame}{Changes in Job Quality and Nature of Work}
AI and automation influence the quality and structure of jobs in several ways:

\begin{itemize}
    \item \textbf{Task Transformation}: Jobs are increasingly divided into tasks, and only certain tasks within jobs are automated.
    \vspace{0.5cm}
    \item \textbf{Increased Productivity}: Automation boosts efficiency and reduces costs, potentially leading to higher economic output and growth.
  \vspace{0.5cm}
    \item \textbf{Work Intensification}: Workers in automated environments may experience greater performance monitoring and pressure, raising concerns about job satisfaction and mental health.
\end{itemize}
\end{frame}

\begin{frame}{Wage Effects}
\begin{itemize}
    \item \textbf{Upward Pressure on High-Skilled Wages}: High-skilled workers who complement automation technologies experience wage gains due to increased productivity and demand.
    \vspace{0.5cm}
    \item \textbf{Downward Pressure on Routine Work}: Workers performing automatable tasks often face wage stagnation or reductions due to reduced demand and increased labor supply competition.
    \vspace{0.5cm}
    \item \textbf{Bargaining Power}: Automation may weaken the bargaining power of workers in heavily automated sectors.
\end{itemize}
\end{frame}

\begin{frame}{Skill Requirements and Education}
Automation drives a shift in the skills demanded by employers:
\vspace{0.5cm}
\begin{itemize}
    \item \textbf{Reskilling and Upskilling}: Workers need to develop skills such as:
    \begin{itemize}
        \item Problem-solving, critical thinking, and creativity.
        \item Digital literacy and technical expertise.
        \item Soft skills like emotional intelligence and communication, particularly for non-automatable jobs.
    \end{itemize}
    \vspace{0.5cm}
    \item \textbf{Lifelong Learning}: Traditional education systems must evolve to emphasize continuous learning.
\end{itemize}
\end{frame}

\begin{frame}{Sectoral Shifts}
The effects of automation vary across industries:
\vspace{0.5cm}
\begin{itemize}
    \item \textbf{Manufacturing and Logistics}: Robotics and AI are transforming production lines, reducing the need for manual labor.
    \vspace{0.5cm}
    \item \textbf{Service Sector}: Automation in customer service and AI in healthcare is augmenting rather than replacing workers.
    \vspace{0.5cm}
    \item \textbf{Creative Industries}: AI is encroaching on creative tasks but human input remains critical for originality.
\end{itemize}
\end{frame}

\begin{frame}{Regional and Demographic Disparities}
The impacts of automation are unevenly distributed:
\vspace{0.5cm}
\begin{itemize}
    \item \textbf{Regional Effects}: Regions dependent on manufacturing face greater job losses, while urban areas are more resilient.
    \vspace{0.5cm}
    \item \textbf{Demographic Effects}: Younger workers may adapt more easily to technological changes, while older workers may face significant challenges in transitioning to new roles.
\end{itemize}
\end{frame}

\begin{frame}{Policy Implications}
Labor economists emphasize the need for proactive policies:
\vspace{0.5cm}
\begin{itemize}
    \item \textbf{Active Labor Market Policies (ALMPs)}: Subsidies for retraining programs, wage insurance, job placement services.
    \vspace{0.5cm}
    \item \textbf{Social Protection}: Strengthening unemployment insurance and exploring Universal Basic Income (UBI).
\vspace{0.5cm}
    \item \textbf{Regulating Automation}: Encouraging responsible adoption of AI and ensuring ethical AI usage to avoid biases in hiring.
\end{itemize}
\end{frame}

\begin{frame}{Long-Term Perspectives}

 \textbf{Economic Growth}: Automation could lead to productivity gains, economic expansion, and higher standards of living.
    \vspace{0.5cm}\\
  \textbf{Redefining Work}: A cultural shift may redefine the role of work in society, focusing on creative and fulfilling roles.

  \vspace{0.5cm}
\textbf{However, }
\begin{itemize}
\item the trajectory of automation’s effects depends on technological advancement, worker adaptability, and policy interventions. 
\item a balanced approach is necessary to ensure inclusive and equitable labor market outcomes.
\end{itemize}
\end{frame}






\begin{frame}{Problem 1: Labor Market Polarization}
\textbf{Scenario:}
Automation has affected 3 categories of workers in a country:
\begin{itemize}
    \item High-skilled workers have seen a 15\% increase in demand.
    \item Low-skilled workers have experienced no change in demand.
    \item Middle-skilled workers have seen a 25\% decrease in demand.
\end{itemize}

Assume that the total labor force consists of 100,000 workers:
\begin{itemize}
    \item 20\% are high-skilled,
    \item 50\% are middle-skilled,
    \item 30\% are low-skilled.
\end{itemize}

\textbf{Question:}
\begin{itemize}
    \item Calculate the total change in labor demand in terms of workers for each skill group.
    \item What is the net effect on the labor market's composition?
\end{itemize}
\end{frame}


\begin{frame}{Problem 1: Labor Market Polarization - Solution}
\textbf{Given:}
\begin{itemize}
    \item Total workforce = 100,000
    \item High-skilled workers = 20\% of 100,000 = 20,000
    \item Middle-skilled workers = 50\% of 100,000 = 50,000
    \item Low-skilled workers = 30\% of 100,000 = 30,000
    \item Changes in demand:
    \begin{itemize}
        \item High-skilled: +15\%
        \item Middle-skilled: -25\%
        \item Low-skilled: 0\%
    \end{itemize}
\end{itemize}


\end{frame}




\begin{frame}{Problem 1: Labor Market Polarization - Solution}


\textbf{Solution:}
\begin{itemize}
    \item High-skilled change in demand: \( 20,000 \times 0.15 = 3,000 \)
    \item Middle-skilled change in demand: \( 50,000 \times (-0.25) = -12,500 \)
    \item Low-skilled change in demand: \( 30,000 \times 0 = 0 \)
    \item High-skilled workers after change: \( 20,000 + 3,000 = 23,000 \)
    \item Middle-skilled workers after change: \( 50,000 - 12,500 = 37,500 \)
    \item Low-skilled workers after change: \( 30,000 \)
\end{itemize}

\textbf{Net Effect:}
\begin{itemize}
    \item High-skilled workers increased by 3,000.
    \item Middle-skilled workers decreased by 12,500.
    \item Low-skilled workers remained unchanged.
\end{itemize}
\end{frame}


\begin{frame}{Problem 2: Job Displacement and Job Creation}
\textbf{Scenario:}
A country has a total workforce of 10 million workers. A study suggests that 40\% of current jobs are highly susceptible to automation, but for every job lost, 0.75 new jobs are created in complementary sectors like AI, data science, and renewable energy.
\vspace{.5cm}\\
\textbf{Question:}
\begin{itemize}
    \item How many jobs are expected to be lost due to automation?
    \item How many new jobs are expected to be created?
    \item What is the net effect on employment in the country? Will the country experience a net increase or decrease in jobs?
\end{itemize}
\end{frame}


\begin{frame}{Problem 2: Job Displacement and Job Creation - Solution}
\textbf{Given:}
\begin{itemize}
    \item Total workforce = 10 million workers
    \item 40\% of jobs are susceptible to automation
    \item For every job lost, 0.75 new jobs are created.
\end{itemize}

\textbf{Solution:}
\begin{itemize}
    \item Jobs displaced = \( 10,000,000 \times 0.40 = 4,000,000 \)
    \item New jobs created = \( 4,000,000 \times 0.75 = 3,000,000 \)
    \item Net job change = \( 3,000,000 - 4,000,000 = -1,000,000 \)
\end{itemize}

\textbf{Conclusion:}
\begin{itemize}
    \item The country will experience a net decrease of 1 million jobs in the short term.
\end{itemize}
\end{frame}







\begin{frame}{Problem 3: Wage Effects}
\textbf{Scenario:}
Consider a country with two sectors:
\begin{itemize}
    \item The \textit{High-Skilled Sector}: The average wage is \$80,000. Automation increases worker productivity by 20\%, leading to a 10\% wage increase for high-skilled workers.
    \item The \textit{Routine Work Sector}: The average wage is \$40,000. Automation reduces the demand for workers, leading to a 5\% decrease in wages.
\end{itemize}

If there are 3 million workers in each sector, calculate the total wage change in each sector.
\vspace{.5cm}\\
\textbf{Question:}
\begin{itemize}
    \item What is the overall effect on total wage payments across both sectors?
\end{itemize}
\end{frame}


\begin{frame}{Problem 3: Wage Effects - Solution}
\textbf{Given:}
\begin{itemize}
    \item High-Skilled Sector:
    \begin{itemize}
        \item Average wage = \$80,000
        \item Productivity increases by 20\%, wage increases by 10\%
        \item Number of workers = 3 million
    \end{itemize}
    \item Routine Work Sector:
    \begin{itemize}
        \item Average wage = \$40,000
        \item Wage decreases by 5\%
        \item Number of workers = 3 million
    \end{itemize}
\end{itemize}

\end{frame}


\begin{frame}{Problem 3: Wage Effects - Solution}


\textbf{Solution:}
\begin{itemize}
    \item High-Skilled Sector:
    \begin{itemize}
        \item Wage increase = \( 80,000 \times 0.10 = 8,000 \)
        \item New wage = \( 80,000 + 8,000 = 88,000 \)
        \item Total wage payment = \( 88,000 \times 3,000,000 = 264,000,000,000 \)
    \end{itemize}
    \item Routine Work Sector:
    \begin{itemize}
        \item Wage decrease = \( 40,000 \times 0.05 = 2,000 \)
        \item New wage = \( 40,000 - 2,000 = 38,000 \)
        \item Total wage payment = \( 38,000 \times 3,000,000 = 114,000,000,000 \)
    \end{itemize}
    \item Total wage payment = \( 264,000,000,000 + 114,000,000,000 = 378,000,000,000 \)
    \item Total wage increase = \( 378,000,000,000 - 360,000,000,000 = 18,000,000,000 \)
\end{itemize}

\textbf{Conclusion:}
\begin{itemize}
    \item The total wage payment increases by \$18 billion across both sectors.
\end{itemize}
\end{frame}




\begin{frame}{Problem 4: Skill Requirements and Education}
\textbf{Scenario:}
A country plans to invest in retraining programs for 1 million displaced workers due to automation. The cost of retraining per worker is \$2,000. After retraining, 70\% of workers are expected to find new jobs at a wage of \$50,000 per year, while the remaining 30\% remain unemployed.
\vspace{.5cm}\\
\textbf{Question:}
\begin{itemize}
    \item What is the total cost of the retraining program?
    \item What is the expected annual increase in wage payments for the retrained workers compared to their previous wage of \$30,000?
\end{itemize}
\end{frame}

\begin{frame}{Problem 4: Skill Requirements and Education - Solution}
\textbf{Given:}
\begin{itemize}
    \item Total displaced workers = 1 million
    \item Cost of retraining per worker = \$2,000
    \item 70\% of retrained workers will find new jobs at \$50,000 per year
    \item 30\% will remain unemployed
    \item Previous wage of displaced workers = \$30,000
\end{itemize}


\end{frame}


\begin{frame}{Problem 4: Skill Requirements and Education - Solution}


\textbf{Solution:}
\begin{itemize}
    \item Total cost of retraining program = \( 1,000,000 \times 2,000 = 2,000,000,000 \) (i.e., \$2 billion)
    \item Number of retrained workers finding new jobs = \( 1,000,000 \times 0.70 = 700,000 \)
    \item Wage increase = \( 50,000 - 30,000 = 20,000 \)
    \item Total annual wage increase = \( 700,000 \times 20,000 = 14,000,000,000 \)
\end{itemize}

\textbf{Conclusion:}
\begin{itemize}
    \item The total cost of the retraining program is \$2 billion.
    \item The expected annual increase in wage payments for retrained workers is \$14 billion.
\end{itemize}
\end{frame}



\begin{frame}{Problem 5: Regional and Demographic Disparities}
\textbf{Scenario:}
The automation of manufacturing has led to job displacement in two regions:
\begin{itemize}
    \item Region A (heavy reliance on manufacturing): 40\% of workers are displaced due to automation.
    \item Region B (diversified economy): 20\% of workers are displaced.
\end{itemize}

The total workforce in Region A is 3 million workers, and in Region B is 5 million workers. Additionally, 30\% of displaced workers are aged 50 and above, while the remaining are under 50.


\end{frame}




\begin{frame}{Problem 5: Regional and Demographic Disparities}


\textbf{Question:}
\begin{itemize}
    \item Calculate the total number of displaced workers in each region.
    \item In Region A, if only 50\% of displaced workers under 50 years old find new employment, how many workers will remain unemployed in Region A?
    \item What percentage of displaced workers in both regions are aged 50 and above?
\end{itemize}
\end{frame}

\begin{frame}{Problem 5: Regional and Demographic Disparities - Solution}
\textbf{Given:}
\begin{itemize}
    \item Region A: Total workforce = 3 million, 40\% displaced.
    \item Region B: Total workforce = 5 million, 20\% displaced.
    \item 30\% of displaced workers are aged 50 and above, while the remaining are under 50.
\end{itemize}


\end{frame}


\begin{frame}{Problem 5: Regional and Demographic Disparities - Solution}

\textbf{Solution:}
\begin{itemize}
    \item Total displaced workers:
    \begin{itemize}
        \item Region A: \( 3,000,000 \times 0.40 = 1,200,000 \)
        \item Region B: \( 5,000,000 \times 0.20 = 1,000,000 \)
    \end{itemize}
    \item Displaced workers aged 50 and above:
    \begin{itemize}
        \item Region A: \( 1,200,000 \times 0.30 = 360,000 \)
        \item Region B: \( 1,000,000 \times 0.30 = 300,000 \)
    \end{itemize}
    \item Displaced workers aged under 50:
    \begin{itemize}
        \item Region A: \( 1,200,000 - 360,000 = 840,000 \)
        \item Region B: \( 1,000,000 - 300,000 = 700,000 \)
    \end{itemize}
    \item Region A unemployment (under 50): \( 840,000 \times 0.50 = 420,000 \)
\end{itemize}


\end{frame}



\begin{frame}{Problem 5: Regional and Demographic Disparities - Solution}

\textbf{Conclusion:}
\begin{itemize}
    \item Total displaced workers: 1,200,000 in Region A and 1,000,000 in Region B.
    \item Displaced workers aged 50 and above: 360,000 in Region A and 300,000 in Region B.
    \item Unemployed in Region A (under 50): 420,000 workers.
\end{itemize}
\end{frame}

\end{document}
